\documentclass[11pt]{article}
\usepackage[T1]{fontenc}
\usepackage{lmodern,microtype}
\usepackage[margin=1in]{geometry}
\usepackage{amsmath,amssymb,amsthm,mathtools}
\usepackage[hidelinks]{hyperref}
\hypersetup{pdftitle={Unbounded normalized cofactor spectra for rank two Gram matrices}}
\newtheorem{theorem}{Theorem}
\newtheorem{lemma}[theorem]{Lemma}
\newtheorem{corollary}[theorem]{Corollary}
\DeclareMathOperator{\per}{per}
\DeclareMathOperator{\rank}{rank}
\DeclareMathOperator{\Ree}{Re}
\newcommand{\CP}{\mathbb{CP}}
\newcommand{\one}{\mathbf 1}
\title{Unbounded normalized cofactor spectra\\for rank two Gram matrices}
\author{Yongxian Zhang\thanks{School of Computer Science and Engineering,
South China University of Technology, Guangzhou, China.
Email: \href{mailto:mxymmxym1@gmail.com}{mxymmxym1@gmail.com}.
ORCID: \href{https://orcid.org/0009-0000-3864-3536}{0009-0000-3864-3536}.}}
\date{8 October 2026}
\hypersetup{pdfauthor={Yongxian Zhang}}
\setlength{\emergencystretch}{3em}
\begin{document}
\maketitle
\begin{abstract}
For a positive-semidefinite Hermitian matrix $A$, let
$C(A)_{ij}=a_{ij}\per A(i\mid j)$, where the minor deletes row $i$ and column $j$.
We prove that $\lambda_{\max}(C(A))/\per A$ has no dimension-independent upper bound,
already for complex rank-two correlation matrices. The same construction makes
$\lambda_{\max}(\Ree C(A))/\per A$ arbitrarily large. The proof computes a spectral
limit under repeated Gram vectors and combines it with projective equidistribution.
It also gives positive-definite examples and unbounded local second-order relative
coefficients for entrywise-positive real correlation perturbations.
\end{abstract}

\section{Statement and scope}
Our convention is that $A(i\mid j)$ deletes row $i$ and column $j$, without
transposing this minor. All Hermitian inner products below are linear in their
first argument. The matrix $C(A)$ is positive semidefinite and
$C(A)\one=(\per A)\one$. Drury disproved the assertion that this prescribed
eigenvalue is always the largest one \cite{drury}. Here we consider all constant
multiples of the permanent.

\begin{theorem}\label{thm:main}
For every $R>0$ there are an integer $N$ and a complex Hermitian correlation
matrix $A$ of rank exactly two such that
\[
 \lambda_{\max}(C(A))>R\per A,
 \qquad \lambda_{\max}(\Ree C(A))>R\per A.
\]
There are also positive-definite correlation matrices with both inequalities.
The positive-definite assertion does not retain rank two.
\end{theorem}

The real part is entrywise, so its largest eigenvalue is the maximum Rayleigh
quotient over real unit vectors. The theorem is an unboundedness assertion as
the order varies; no rate in $N$ is asserted. It is distinct from unbounded
Hadamard-product bunching ratios \cite{seron}, and from bounds for the full
Schur-power matrix. It does not resolve Lieb's normalized character or subgroup
permanental-dominance conjecture \cite{wanless}. A separate source audit records
the finite literature search; no claim of exhaustive priority verification is made.

\section{A repetition limit}
Let $v_i=(a_i,b_i)$ be nonzero unit rows, $1\le i\le n$, put
$\ell_i(x,y)=a_i x+b_i y$, and set $P=\prod_i\ell_i$ and $A_{ij}=v_i v_j^*$.
For binary forms of degree $d$ define
\[
 \left\langle\sum_{k=0}^d p_kx^{d-k}y^k,
 \sum_{k=0}^d q_kx^{d-k}y^k\right\rangle_B
 =\sum_{k=0}^d(d-k)!k!p_k\overline{q_k}.
\]
Leibniz expansion gives
\begin{equation}\label{eq:gram}
 \left\langle\prod_i\ell_i,\prod_j\ell'_j\right\rangle_B
 =\per[\langle\ell_i,\ell'_j\rangle_B].
\end{equation}
In particular $\per A=\|P\|_B^2>0$, and $C(A)$ is the Gram matrix of
$v_i\otimes(P/\ell_i)$. This proves its positivity with the stated conventions.
For Haar probability measure $\mu$ on $\CP^1$, represented by unit vectors,
\begin{equation}\label{eq:sphere}
 \langle p,q\rangle_B=(d+1)!\int p\overline q\,d\mu.
\end{equation}
Indeed, phase integration makes distinct monomials orthogonal and
$\int|x|^{2a}|y|^{2b}\,d\mu=a!b!/(a+b+1)!$.

Repeat each row $L$ times, write $N=nL$, and let $A_L$ be the resulting Gram
matrix. Then $P_L=P^L$ and, with cluster indices $(i,a)$,
\begin{equation}\label{eq:cluster}
 C(A_L)_{(i,a),(j,b)}
 =A_{ij}\langle P^L/\ell_i,P^L/\ell_j\rangle_B.
\end{equation}
Let $U_L:\mathbb C^n\to\mathbb C^{nL}$ map $e_i$ to the vector equal to
$L^{-1/2}$ on cluster $i$ and zero elsewhere. All clusters in \eqref{eq:cluster}
are constant, so $C(A_L)$ vanishes on $\operatorname{ran}(U_L)^\perp$.
Its nonzero normalized spectrum, including multiplicities, is that of
$K_L=U_L^*C(A_L)U_L/\per A_L$.
Writing $d\mu_L=|P|^{2L}d\mu/\int|P|^{2L}d\mu$, the factorials in
\eqref{eq:sphere} give the exact formula
\begin{equation}\label{eq:compression}
 (K_L)_{ij}=\frac{L}{nL+1}A_{ij}
       \int\frac{d\mu_L}{\ell_i\overline{\ell_j}}.
\end{equation}
The expression against the density at a zero is interpreted through the
polynomial $(P^L/\ell_i)\overline{(P^L/\ell_j)}$.

\begin{lemma}\label{lem:limit}
Suppose $|P|$ has a unique projective maximum. Rotate it to $(1,0)$ and put
$z_i=b_i/a_i$. Then every $a_i\ne0$, $\sum_i z_i=0$, and
\[
 K_L\longrightarrow\frac{J_n+zz^*}{n}
 \quad\hbox{in operator norm}.
\]
Consequently
\begin{align}
 \lim_{L\to\infty}\frac{\lambda_{\max}(C(A_L))}{\per A_L}
 &=\max\left(1,\frac{\|z\|^2}{n}\right),\label{eq:speclimit}\\
 \lim_{L\to\infty}\frac{\lambda_{\max}(\Ree C(A_L))}{\per A_L}
 &\ge\max\left(1,\frac{\|z\|^2}{2n}\right).\label{eq:reallimit}
\end{align}
\end{lemma}
\begin{proof}
The maximum of the nonzero polynomial is positive, so no factor vanishes there.
The integrand in \eqref{eq:compression} is continuous near the maximum $u$.
Outside any neighborhood of $u$, $|P|\le a<|P(u)|$. On a smaller neighborhood
of positive measure, $|P|\ge b>a$. The absolute unnormalized tail is bounded by
\[
 a^{2L-2}\sup\left|(P/\ell_i)(P/\ell_j)\right|,
\]
whereas the denominator is bounded below by a positive constant times $b^{2L}$.
Thus the apparent poles cause no problem and the integral tends to
$(a_i\overline{a_j})^{-1}$. Since $n$ is fixed, entrywise convergence implies
operator-norm convergence and gives the stated matrix limit.

In the local coordinate $(1,t)/\sqrt{1+|t|^2}$,
\[
 \log|P|^2=\text{constant}+\sum_i\log|1+z_i t|^2
                    -n\log(1+|t|^2).
\]
The two real first derivatives at zero vanish, hence $\sum_i z_i=0$.
The limiting matrix therefore has eigenvalues $1$, $\|z\|^2/n$, and zeros,
which proves \eqref{eq:speclimit}. For $z=x+iy$ its real part is
$(J_n+xx^T+yy^T)/n$, with $x,y\perp\one$. The last two terms have rank at most
two and trace $\|z\|^2/n$, proving \eqref{eq:reallimit}. Compression commutes
with real part because $U_L$ is real.
\end{proof}

\section{Equidistribution and selection of one maximum}
Choose collections of $m$ unit rows whose empirical projective measures
converge weakly to Haar measure on $\CP^1$. Such a sequence is obtained from
midpoint grids in $p=|a|^2\in(0,1)$ and the relative phase, with both grid sizes
tending to infinity. Let $u_m$ be any maximum of the corresponding product,
and rotate variables so that $u_m=(1,0)$.

These moving rotations preserve weak convergence to Haar. Indeed, every
subsequence of rotations has a convergent subsubsequence by compactness of
$U(2)$. For each continuous test function, the rotated functions then converge
uniformly; weak convergence and Haar invariance give the same limiting integral.
This subsequence argument proves the claim for the full sequence.

For $T>0$ the function
\[
 f_T([a:b])=\min(|b/a|^2,T),\qquad f_T([0:1])=T,
\]
is continuous. Since $|a|^2$ is uniform on $[0,1]$ under Haar measure,
\begin{equation}\label{eq:truncation}
 \int f_T\,d\mu=\int_0^1\min\left(\frac{1-p}{p},T\right)dp
 =\log(1+T).
\end{equation}
Every $a_i$ in the rotated configuration is nonzero. It follows that
\[
 \liminf_{m\to\infty}\frac1m\sum_{i=1}^m|b_i/a_i|^2\ge\log(1+T)
 \quad\text{for every }T>0.
\]
Hence these averages tend to infinity. Only bounded continuous test functions
have been used; there is no assumption of convergence of an unbounded moment.

To make the maximum unique, append one unit linear form $\ell_u(z)=\langle z,u_m\rangle$,
linear in $z$. On the unit sphere its modulus is at most one, with equality
only at the projective point $u_m$. Multiplication by this factor leaves the
old maximum value unchanged and makes its location unique. In the rotated
coordinates the added row is $(1,0)$, contributing $z=0$. For the augmented
base $n=m+1$ we still have
\begin{equation}\label{eq:divergence}
 \frac1n\sum_{i=1}^n|z_i|^2
 =\frac1{m+1}\sum_{i=1}^m|b_i/a_i|^2\longrightarrow\infty.
\end{equation}

\begin{proof}[Proof of Theorem \ref{thm:main}]
Given $R$, first choose one finite base in \eqref{eq:divergence} with average
exceeding $2R+2$. Equidistribution ensures that a sufficiently large base spans
$\mathbb C^2$. Freeze it and then take a sufficiently large integer $L$ in
Lemma~\ref{lem:limit}. Both ratios exceed $R$, all rows are unit, and the Gram
matrix has rank exactly two.

For a strict witness $A$, polynomial dependence of $C(A)$ and $\per A$, continuity
of Hermitian eigenvalues, and $\per A>0$ show that the strict inequalities persist
for $A+\varepsilon I$ with sufficiently small $\varepsilon>0$. The matrix
$(A+\varepsilon I)/(1+\varepsilon)$ is a positive-definite correlation matrix.
Both ratios are unchanged by this final scalar rescaling, since
$C(cA)=c^N C(A)$ and $\per(cA)=c^N\per A$.
\end{proof}

\section{A local consequence and an exact finite certificate}
The following is the real-direction specialization of the perturbation formula
of Pioge and coauthors \cite[Lemma~1]{pioge}; we include the calculation to fix
the Gram and conjugation conventions. For a real unit vector $w$, let
\[
 B(\varepsilon)_{ij}=
 \frac{1+\varepsilon^2w_iw_j}
 {\sqrt{(1+\varepsilon^2w_i^2)(1+\varepsilon^2w_j^2)}}.
\]
This is a real correlation matrix of rank at most two, with strictly positive
entries for sufficiently small $\varepsilon$. Expanding the entries and then
the permanent gives
\begin{equation}\label{eq:local}
 \frac{\per(A\circ B(\varepsilon))}{\per A}
 =1+\varepsilon^2\left(\frac{w^T\Ree C(A)w}{\per A}-1\right)
       +O(\varepsilon^4).
\end{equation}
In fact the quadratic coefficient before division by $\per A$ is
\[
 \sum_{ij}C(A)_{ij}\left(w_iw_j-\frac{w_i^2+w_j^2}{2}\right)
 =w^TC(A)w-\per A=w^T\Ree C(A)w-\per A,
\]
using both row and column sums. Theorem~\ref{thm:main} therefore makes the
relative coefficient in \eqref{eq:local} arbitrarily large. This is a local
statement; it gives no uniform finite-$\varepsilon$ gain or perturbative interval.
For complex $w$, the column-Gram convention in \cite{pioge} has numerator
$1+\varepsilon^2\overline{w_i}w_j$; replacing it by the opposite convention
would conjugate the Rayleigh direction.

For an optional finite check, the companion file \texttt{witness\_n200.csv}
specifies $200$ Gaussian-integer rows $v_i\in\mathbb Z[i]^2$ and a vector
$w\in\mathbb Z[i]^{200}$. The standard-library integer verifier proves
\[
 \frac{269}{100}<\frac{w^*C(VV^*)w}{\per(VV^*)\|w\|^2}<\frac{27}{10}.
\]
It independently forms each omitted-factor polynomial and evaluates
$\sum_{a=1}^2\|\sum_i\overline{w_i}v_{ia}(P/\ell_i)\|_B^2$.
This certificate illustrates a finite violation and checks the conventions;
unboundedness follows from the analytic proof, not from numerical extrapolation.


\section*{Source edition and reproduction}
This signed article edition retains the mathematical argument of the
previously archived source: \url{https://doi.org/10.5281/zenodo.23252140}.
The original source, exact inputs, checking programs and prior review
records remain available there. The companion README identifies the
frozen Git version and the scope of each auxiliary computation.
In a writable copy of that archived source directory, run:
\begin{verbatim}
python3 verify_witness.py
python3 verify_witness_independently.py
\end{verbatim}
The computational checks do not replace the analytic proof.

\section*{Funding and research methods}
This research was conducted independently and received no external
funding. OpenAI AI assistants assisted mathematical research, verification
code and internal review. Mathematical claims are supported by the
written proof and the specified public evidence; AI output is not itself
proof evidence. Internal AI-assisted review is not external human peer
review.
\section*{Declaration of generative AI and AI-assisted technologies
in the writing process}
OpenAI AI assistants assisted drafting, organization and LaTeX editing.
The author takes responsibility for the manuscript and its descriptions
of the proof and verification evidence.

\begin{thebibliography}{9}
\bibitem{drury} S.~Drury, A counterexample to a question of Bapat \& Sunder,
\emph{Math. Inequal. Appl.} \textbf{21} (2018), 517--520.
\href{https://doi.org/10.7153/mia-2018-21-37}{doi:10.7153/mia-2018-21-37}.
\bibitem{pioge} L.~Pioge, K.~K.~Pietrasz, B.~Seron, L.~Novo and N.~J.~Cerf,
A logical implication between two conjectures on matrix permanents,
\emph{Linear Algebra Appl.} \textbf{725} (2025), 309--318.
\href{https://doi.org/10.1016/j.laa.2025.07.011}{doi:10.1016/j.laa.2025.07.011};
\href{https://arxiv.org/abs/2508.00111}{arXiv:2508.00111}.
\bibitem{seron} B.~Seron, L.~Novo and N.~J.~Cerf,
Boson bunching is not maximized by indistinguishable particles,
\emph{Nature Photonics} \textbf{17} (2023), 702--709.
\href{https://doi.org/10.1038/s41566-023-01213-0}{doi:10.1038/s41566-023-01213-0}.
\bibitem{wanless} I.~M.~Wanless, Lieb's permanental dominance conjecture,
In \emph{The Physics and Mathematics of Elliott Lieb}, vol.~II,
EMS Press (2022), 501--516.
\href{https://doi.org/10.4171/90-2/48}{doi:10.4171/90-2/48};
\href{https://arxiv.org/abs/2202.01867}{arXiv:2202.01867}.
\end{thebibliography}
\end{document}
